%LTeX: language=es-AR

\documentclass[12pt]{amsart}

%Fuente
\usepackage[utf8]{inputenc}
\usepackage[spanish]{babel}
\usepackage[T1]{fontenc}
\usepackage{lmodern}
\usepackage{epigraph}
\usepackage{lipsum}
\usepackage[a4paper, margin=2.5cm]{geometry}
\usepackage{multicol}
\usepackage{amsmath}
\usepackage{amssymb}
\usepackage[shortlabels]{enumitem}
\usepackage{listings}
\lstdefinestyle{micode}{
    basicstyle=\ttfamily\small,
    breaklines=true,
    columns=fullflexible,
    keepspaces=true,
    showstringspaces=false
}

\DeclareMathOperator{\mcd}{mcd}

\renewcommand{\refname}{Bibliografía}

%diagramas
\usepackage{tikz-cd}

%enumerar páginas y encabezados
\pagestyle{headings}

%Comandos para escribir C más rápido
\usepackage{mathrsfs}
\newcommand{\CC}{\mathbb{C}}
\newcommand{\NN}{\mathbb{N}}
\newcommand{\QQ}{\mathbb{Q}}
\newcommand{\RR}{\mathbb{R}}
\newcommand{\ZZ}{\mathbb{Z}}
\newcommand{\FF}{\mathbb{F}}
\newcommand{\cont}{\mathcal{C}}

\newcommand{\pp}{\mathfrak{p}}

\newcommand{\aaa}{\mathfrak{a}}
\newcommand{\bbb}{\mathfrak{b}}

\newcommand{\OO}{\mathcal{O}}

\newcommand{\leg}[2]{\left( \frac{#1}{#2} \right)}
\newcommand{\minimat}[4]{\left(\begin{smallmatrix} #1 & #2 \\ #3 & #4 
\end{smallmatrix}\right)}
\newcommand{\lp}{\left(}
\newcommand{\rp}{\right)}
\newcommand{\lb}{\left|}
\newcommand{\rb}{\right|}
\newcommand{\lc}{\left<}
\newcommand{\rc}{\right>}
\newcommand{\lco}{\left[}
\newcommand{\rco}{\right]}

%dobles implicaciones
\usepackage{csquotes}

%Entornos con estilo 'tipo teorema'
\theoremstyle{plain}
\newtheorem{teo}[equation]{Teorema}
\newtheorem{ej}[equation]{Ejercicio}
\newtheorem{prop}[equation]{Proposición}
\newtheorem{lem}[equation]{Lema}
\newtheorem{coro}[equation]{Corolario}
\newtheorem{defi}[equation]{Definición}
\newtheorem{obs}[equation]{Observación}
\newtheorem*{nota*}{Notación}
\renewcommand*{\proofname}{Demostración}

% Símbolo de Tate--Shafarevich
\DeclareFontFamily{U}{wncy}{}
\DeclareFontShape{U}{wncy}{m}{n}{<->wncyr10}{}
\DeclareSymbolFont{cyrletters}{U}{wncy}{m}{n}
\DeclareMathSymbol{\Sha}{\mathord}{cyrletters}{"58}

%Hipervínculos cliqueables, en color azul
\usepackage[colorlinks=true,allcolors=blue]{hyperref}

\title{Título}

\author{Pedro Nicolás Peña}

\begin{document}

\maketitle

% con Alt + Z te acomoda las lineas

\section*{Sección}

Vamos a trabajar mucho con la curva $E_t: y^2 = x(x+1)(x+t^2) = x^3+(t^2+1)x^2 + t^2x$ porque tiene grupo de torsión isomorfo a $\ZZ/2\ZZ \oplus \ZZ/4\ZZ$ sobre $\QQ$ y si $t^2-1$ es un cuadrado entonces sobre $\QQ(i)$ es $\ZZ/4\ZZ \oplus \ZZ/4\ZZ$. Otra presentación de la curva es $E_{mn}: y^2 = x(x+n^2)(x+m^2)$ con la condición siendo que $m^2-n^2$ sea un cuadrado. Y más en general la curva va a ser $E_{rs}: y^2 = x(x+r^2)(x+s^2)$ donde $r$ y $s$ van a ser polinomios en las variables $t$ o $m$ y $n$. El cambio de variables es fácil de ver: si $t = \frac{r}{s}$ entonces

\[
y^2 = x(x+1)\lp x+\leg{r}{s}^2\rp \iff {(y's^{-3})}^2 = (x's^{-2})((x's^{-2})+1)\lp (x's^{-2})+\leg{r}{s}^2\rp 
\]

\[
\iff y'^2 s^{-6} = x'(x'+s^2)(x'+r^2) s^{-6} \iff y'^2 = x'(x'+r^2)(x'+s^2)
\]


Notemos también que es la curva de Legendre $E_\lambda : y^2 = x(x+1)(x+\lambda)$ con $\lambda = t^2$, sobre la cual hay mucha literatura. Importante notar que la superficie elíptica (de Legendre) tiene rango 0, osea no tiene puntos independientes para toda especializacion.

\textbf{La  curva isógena:}

La isogenia con kernel generado por $(0,0)$ va a la curva $E'_t: y^2 = x^3-2(t^2+1)x^2+{(t^2-1)}^2x$ o en terminos de $r$ y $s$ va a $E'_{rs}: y^2 = x^3-2(r^2+s^2)x^2+{(r^2-s^2)}^2x$.



\textbf{Los puntos:}

Sí ponemos $P = (rs,rs(r+s))$ y $Q = (r(u-r), iru(u-r))$ con $u^2=r^2-s^2$

\[
\begin{array}{|c|c|c|c|c|}
\hline
+   & \OO    & P      & 2*P    & 3*P \\ \hline

\OO & \OO     & (rs,rs(r+s)) & (0,0)     & (rs,-rs(r+s)) \\ \hline

Q& (-r(r-u),  & (-s(s-iu), & (-r(r+u), & (-s(s+iu), \\

&  -iru(r-u)) & su(s-iu))  & iru(r+u)) & -su(s+iu)) \\ \hline

2*Q& (-r^2,0) & (-rs,-rs(r-s))& (-s^2,0)& (-rs,rs(r-s)) \\ \hline

3*Q & (-r(r-u),& (-s(s+iu, & (-r(r+u), & (-s(s-iu), \\ 

&   iru(r-u)) & su(s+iu)) & -iru(r+u)) & -su(s-iu)) \\ \hline
\end{array}
\]

(esta tabla está chequeada a mano con Sage)

\textbf{Discriminante:}

\[
\Delta_{Ers} = 16 r^4 s^4 {(r^2-s^2)}^2 = 2^4 r^4 s^4 u^4
\]


\textbf{j-invariante:}

El numerador tiene pinta de que pasa por todos los valores.

\[
j(E_{rs}) = 2^8 \frac{{(r^4 - r^2s^2 + s^4)}^3}{r^4s^4{(r^2-s^2)}^2}
\]


\textbf{Selmer:}

Tenemos que para cada $d \in \QQ^\times / {(\QQ^\times)}^2$ la curva $\cont_d$ tiene soluciones en $\QQ_v$ para toda completación de $\QQ$ si y solo si $d$ está en el Selmer. Notar que la curva $\cont_d$ usa los coeficientes de la $E'$ (y $\cont_d'$ usará los de E). Solamente tenemos que chequear los $d$ para los que dividen a $b'$.

\[
\cont_d : y^2 = dx^4 -2(t^2+1) x^2+ {(t^2-1)}^2/d
\]

\[
\cont_d' : y^2 = dx^4 +(t^2+1) x^2+ t^2/d
\]

\textbf{Mas puntos:}

Sí queremos más puntos necesitamos que $x(x+r^2)(x+s^2)$ sea un cuadrado en el cuerpo (recordar que estamos en $\QQ(i)$), por ejemplo podemos pedir $\pm x = (x+r^2)(x+s^2)$ es decir $0 = x^2+(r^2+s^2\pm 1)x +r^2s^2$ es decir que ${(r^2+s^2\pm 1)}^2-4r^2s^2$ sea un cuadrado. Sí ponemos el resto de fórmulas queremos que los siguientes sean cuadrados en el cuerpo:
\[
{(r^2+s^2\pm 1)}^2-4r^2s^2 
\]
\[
{(r^2\pm 1)}^2\mp 4s^2
\]
\[
{(s^2\pm 1)}^2\mp 4r^2
\]

Obvio habría que ver que estos puntos no terminen siendo los de torsión. 

Tal vez el teorema que me dé puntos independientes no sea estilo Green-Tao sino con ternas pitagóricas...

Ejemplo (me olvidé de la condición de que $r^2-s^2$ sea un cuadrado):

\begin{lstlisting}[style=micode]
Qi.<i> = QuadraticField(-1)

E0 = EllipticCurve(QQ, [0, 7^2 + 24^2, 0, 7^2*24^2, 0])
Ee = E0.change_ring(Qi)   #mejor esto que definirla directo en Qi pq asi tenemos .rank() y demas

P = Ee(-18, i*558)
print(Ee.rank(known_points=[P])) => 3
print(Ee.gens()) => ((-18 : 558*i : 1), (504*i - 504 : 1512*i - 15624 : 1), (-525 : 3570 : 1))

Q = Ee(-32, i*544) 
print(2*P == -2*Q) => True
print(P == Ee(-24^2,0) - Q) => True
\end{lstlisting}



\begin{prop}
Un punto $(x_1,y_1)$ en $E(k): y^2 = (x-\alpha_1)(x-\alpha_2)(x-\alpha_3)$ se puede dividir por $2$ si y solo si los números $x_1-\alpha_i$ son cuadrados en $k$.
\end{prop}

\begin{prop}
Para una curva elíptica definida sobre $K$, una extensión cuadrática $L=K(\sqrt{D})$ y el twist $E_D$ vale

\[
rg(E(L)) = rg(E(K)) + rg(E_D(K))
\]
\end{prop}






\textbf{Que se sabe:}

Hay infinitas curvas de rango $r$ para $0\leq r \leq 4$ sobre todo cuerpo de números. Las da Zywina en \cite{ZywinaSmallRank} que todas tienen full 2-torsion (generada por $(0,0)$)(obvio podes matarla y tenes de torsion trivial) con $\Sha [2] = 0$ (falta leer un poco mas de ese). Todo vino de una secuencia de papers: \cite{ZywinaRankNotJump}, \cite{ZywinaRankTwo} y \cite{ZywinaRankOneStability}. Usa numeros de Tamagawa para calcular los simbolos de Kodaira y el global root number porque (la conjetura de paridad dice que) te da la paridad del rango, y estas cosas te dan restricciones sobre las especializaciones.



De Elkies hay curvas con rango alto, todas sobre $\QQ$, y ordenadas por grupo de torsion: Para el trivial el record es $29$ y está en un post random, y para varios grupos ($\ZZ /n\ZZ$ para $n\in \{2,3,4,6,7\}$) hay un paper: \cite{ElkiesKlagsbrun2020}. Lo hace con superficies K3.


Sobre la curva de Legendre:\cite{Ulmer2014}.
















\newpage
\begin{thebibliography}{99}

\bibitem{Ulmer2014}
D. Ulmer, \emph{Explicit points on the Legendre curve},
J. Number Theory \textbf{136} (2014), 165--194.

\bibitem{ZywinaRankNotJump}
D. Zywina, \emph{An elliptic surface with infinitely many fibers for which the rank does not jump},
preprint, \href{https://arxiv.org/abs/2502.01026}{arXiv:2502.01026}, 2025.

\bibitem{ZywinaRankTwo}
D. Zywina, \emph{There are infinitely many elliptic curves over the rationals of rank 2},
preprint, \href{https://arxiv.org/abs/2502.01957}{arXiv:2502.01957}, 2025.

\bibitem{ZywinaRankOneStability}
D. Zywina, \emph{Rank one elliptic curves and rank stability},
preprint, \href{https://arxiv.org/abs/2505.16960}{arXiv:2505.16960}, 2025.

\bibitem{ZywinaSmallRank}
D. Zywina, \emph{On the infinitude of elliptic curves over a number field with prescribed small rank},
preprint, \href{https://arxiv.org/abs/2602.10865}{arXiv:2602.10865}, 2026.

\bibitem{ElkiesKlagsbrun2020}
N. D. Elkies and Z. Klagsbrun, \emph{New rank records for elliptic curves having rational torsion},
preprint, \href{https://arxiv.org/abs/2003.00077}{arXiv:2003.00077}, 2020.

\end{thebibliography}


\end{document}